\documentclass[../../main.tex]{subfiles}
\begin{document}

\chapter{Example}\label{ch:example}

This chapter exists to prove the wiring end to end, and to show the three
conventions the tooling depends on. Delete it once your own first chapter
replaces it.

\section{A figure that comes from a notebook}\label{sec:example_figure}

\Cref{Fig: example_scatter} is not committed by hand. It is written into this
repository by the analysis template's \texttt{save\_fig}, which produces the PNG
below \emph{and} an interactive HTML export from one figure object, so the two can
never disagree about what the data says.

Note what the block does not contain: no path, no extension guessing, no
\texttt{width=} in points. The filename is bare because \verb|\graphicspath| in
\texttt{document\_settings.sty} resolves it, and the width is \verb|\width| so
every figure in the document changes together.

\begin{figure}[H]
    \centering
    \includegraphics[width=\width]{example_scatter.png}
    \caption{Synthetic measurements against their reference values. Generated by
      \texttt{notebooks/example/example\_figure.ipynb} in the analysis template.}
    \label{Fig: example_scatter}
\end{figure}

The label is \verb|Fig: <stem>|, matching the figure's filename. That is the
convention \texttt{save\_fig} uses to decide whether this figure is already
referenced: because this block exists, re-running the notebook will not append a
duplicate. Rename one without the other and you get two blocks for one figure.

\section{A table that comes from a notebook}\label{sec:example_table}

Tables arrive the same way, but through a library file rather than as separate
artefacts. \texttt{save\_table} writes a booktabs tabular into
\texttt{tables.tex}, wrapped in \texttt{catchfilebetweentags} markers, and you
pull the one you want into the prose:

\ExecuteMetaData[Chapters/Example/tables.tex]{example_summary}

Re-saving a table replaces it in place inside that file, so the numbers follow
the analysis without the surrounding prose moving.

\section{Citations and glossary}\label{sec:example_conventions}

A citation looks like \textcite{example2024}, and an acronym expands on first use
as \gls{tbd} and contracts thereafter: \gls{tbd}. Both are ordinary LaTeX; they
are mentioned here only because a publishing agent has to resolve them when it
flattens this document, and \texttt{.doc-publish/macros.py} is where it is told
how.

\end{document}
